\documentclass[12pt,a4paper]{amsart}

\usepackage{amsmath,amssymb,amsthm}
\usepackage{mathtools}
\usepackage{hyperref}
\usepackage{geometry}
\geometry{margin=2.5cm}
\usepackage{booktabs}

% -------------------------------------------------------
\newtheorem{theorem}{Theorem}[section]
\newtheorem{lemma}[theorem]{Lemma}
\newtheorem{corollary}[theorem]{Corollary}
\newtheorem{proposition}[theorem]{Proposition}
\theoremstyle{definition}
\newtheorem{definition}[theorem]{Definition}
\newtheorem{remark}[theorem]{Remark}

% -------------------------------------------------------
\newcommand{\N}{\mathbb{N}}
\newcommand{\Z}{\mathbb{Z}}
\newcommand{\Q}{\mathbb{Q}}
\newcommand{\R}{\mathbb{R}}
\newcommand{\C}{\mathbb{C}}
\newcommand{\eq}[1]{e\!\left(#1\right)}

% -------------------------------------------------------
\begin{document}

\title{The Goldbach Singular Series and the\\
       Hardy--Littlewood Prime $k$-Tuples Conjecture}

\author{Michael Fuhrmann}
\address{Specialist Mathematics Project, Build 80}
\email{specialist-maths@localhost}
\date{2026}

\begin{abstract}
The Hardy--Littlewood conjectures (1923) provide precise quantitative
predictions for additive problems involving primes.  For binary Goldbach,
the predicted number of representations of a large even integer $n$ as
$n = p + q$ (summing over $\log p \cdot \log q$) is
\[
  R_2(n) \;\sim\; \mathfrak{S}(n) \cdot n,
\]
where the \emph{Goldbach singular series}
\[
  \mathfrak{S}(n)
  \;=\; 2 C_2 \prod_{\substack{p \mid n \\ p > 2}} \frac{p-1}{p-2},
  \qquad
  C_2 \;=\; \prod_{p \geq 3} \frac{p(p-2)}{(p-1)^2} \approx 0.6601618\ldots
\]
encodes the local arithmetic of $n$.  We derive this formula, prove
$\mathfrak{S}(n) > 0$ for all even $n \geq 4$, compute $\mathfrak{S}(n)$
for the first few even integers, and explain why the circle method
\emph{cannot} prove the binary Goldbach conjecture unconditionally
(the ``parity problem'').  We also state the Hardy--Littlewood prime
$k$-tuples conjecture~$\mathbf{H}$, which unifies and quantifies all
prime constellation problems, and discuss its relationship to the
twin prime and Goldbach conjectures.
\end{abstract}

\maketitle
\tableofcontents

% ================================================================
\section{Introduction}
% ================================================================

Binary Goldbach's conjecture asserts that every even integer $n \geq 4$
is the sum of two primes.  Numerically verified for $n \leq 4 \times 10^{18}$
(Oliveira e Silva et al., 2014), it remains unproven.  Unlike ternary
Goldbach (proved by Helfgott, Paper~18), binary Goldbach lies beyond the
reach of the circle method alone, due to the parity problem.

However, the \emph{expected} number of representations is well-understood.
Define the \emph{weighted Goldbach representation count}:
\[
  R_2(n) \;=\; \sum_{\substack{p_1 + p_2 = n \\ p_i \leq n}} \log p_1 \cdot \log p_2
  \;=\; \int_0^1 S(\alpha)^2\, \eq{-n\alpha}\, d\alpha,
\]
where $S(\alpha) = \sum_{m \leq n} \Lambda(m) \eq{m\alpha}$ is the prime
exponential sum (Paper~18).

The Hardy--Littlewood circle method predicts:

\begin{conjecture}[Hardy--Littlewood, 1923]\label{conj:HL}
  For every even $n \geq 4$:
  \[
    R_2(n) \;\sim\; \mathfrak{S}(n) \cdot n,
  \]
  where $\mathfrak{S}(n)$ is the Goldbach singular series.  Equivalently,
  the number of representations $n = p + q$ satisfies:
  \[
    r_2(n) \;=\; \#\{p : p \text{ prime},\; n-p \text{ prime}\}
    \;\sim\; \frac{\mathfrak{S}(n) \cdot n}{(\log n)^2}.
  \]
\end{conjecture}

\textbf{Notation.}  $\eq{x} = e^{2\pi ix}$.  $\Lambda$ = von Mangoldt function.
$C_2 = \prod_{p \geq 3} p(p-2)/(p-1)^2 \approx 0.6601618$ = twin prime constant.
$\phi(q)$ = Euler's totient.  $\mu(q)$ = M\"obius function.

% ================================================================
\section{The Goldbach Singular Series}
\label{sec:singular}
% ================================================================

\subsection{Derivation from the circle method}

We derive $\mathfrak{S}(n)$ by formally carrying out the major arc
analysis for $\int_0^1 S(\alpha)^2 \eq{-n\alpha}\, d\alpha$.

By Lemma~3.1 of Paper~18, on the major arc $\mathfrak{M}(a/q)$ with
$\gcd(a,q)=1$ and $q \leq Q$:
\[
  S(\alpha) \;\approx\; \frac{\mu(q)}{\phi(q)}\, I(\beta),
  \quad \alpha = \frac{a}{q} + \beta.
\]
Therefore:
\[
  \int_{\mathfrak{M}(a/q)} S(\alpha)^2\, \eq{-n\alpha}\, d\alpha
  \;\approx\; \left(\frac{\mu(q)}{\phi(q)}\right)^2 \eq{-\frac{an}{q}}
               \int_{-Q/N}^{Q/N} I(\beta)^2\, \eq{-n\beta}\, d\beta.
\]
Summing over $q \leq Q$ and $\gcd(a,q)=1$:
\begin{align*}
  \int_{\mathfrak{M}} S^2\, \eq{-n\cdot}
  &\;\approx\; \left(\sum_{q=1}^{Q} \frac{\mu(q)^2}{\phi(q)^2}
                \sum_{\substack{a=1\\ \gcd(a,q)=1}}^{q} \eq{-\frac{an}{q}}\right)
               \cdot \int_{-\infty}^{\infty} I(\beta)^2\, \eq{-n\beta}\, d\beta \\
  &\;=\; \left(\sum_{q=1}^{\infty} \frac{\mu(q)^2\, c_q(n)}{\phi(q)^2}\right)
          \cdot n \;+\; \text{error},
\end{align*}
where $c_q(n) = \sum_{\gcd(a,q)=1} \eq{an/q}$ is the Ramanujan sum.

\begin{definition}[Goldbach singular series]\label{def:Goldbach_sing}
  \[
    \mathfrak{S}(n)
    \;=\; \sum_{q=1}^{\infty} \frac{\mu(q)^2\, c_q(n)}{\phi(q)^2}.
  \]
  This series converges absolutely: for squarefree $q$ one has
  $|c_q(n)/\phi(q)^2| \leq 1/\phi(q)$, and more precisely
  $|c_q(n)/\phi(q)^2| \leq 1/\phi(q)^2$ when the Ramanujan sum is
  estimated via multiplicativity.  Since
  $\sum_q \mu(q)^2/\phi(q)^2 = \prod_p (1 + 1/(p-1)^2) < \infty$,
  the series converges absolutely.
\end{definition}

\subsection{Euler product representation}

\begin{theorem}[Euler product for $\mathfrak{S}(n)$]\label{thm:euler}
  For even $n \geq 2$:
  \[
    \mathfrak{S}(n)
    \;=\; 2 C_2 \prod_{\substack{p \mid n \\ p \geq 3}} \frac{p-1}{p-2},
  \]
  where $C_2 = \prod_{p \geq 3} \frac{p(p-2)}{(p-1)^2}$ is the
  twin prime constant.
\end{theorem}

\begin{proof}
  The Ramanujan sum $c_q(n)$ is multiplicative in $q$ and equals
  $\mu(q/\gcd(q,n)) \phi(q)/\phi(q/\gcd(q,n))$.  Since $\mu(q)^2 = 1$
  for squarefree $q$ and $0$ otherwise, the sum over $q$ reduces to
  squarefree values.

  For squarefree $q = p$ prime:
  \[
    c_p(n) \;=\; \sum_{\substack{a=1\\ \gcd(a,p)=1}}^{p} \eq{\frac{an}{p}}
    \;=\; \begin{cases} p-1 & \text{if } p \mid n, \\ -1 & \text{if } p \nmid n. \end{cases}
  \]
  For squarefree $q$: $c_q(n) = \prod_{p \mid q} c_p(n)$ (multiplicativity).

  The sum $\mathfrak{S}(n) = \sum_{q \text{ squarefree}} c_q(n)/\phi(q)^2$
  factors over primes as an Euler product:
  \[
    \mathfrak{S}(n)
    \;=\; \prod_{p} \left(1 + \frac{c_p(n)}{\phi(p)^2}\right)
    \;=\; \prod_{p} \left(1 + \frac{c_p(n)}{(p-1)^2}\right).
  \]

  For $p = 2$ (even $n$): $c_2(n) = (-1)^n = 1$, factor = $1 + 1/1^2 = 2$.
  For $p \geq 3$ with $p \mid n$: $c_p(n) = p-1$, factor = $1 + (p-1)/(p-1)^2 = 1 + 1/(p-1)$.
  For $p \geq 3$ with $p \nmid n$: $c_p(n) = -1$, factor = $1 - 1/(p-1)^2 = ((p-1)^2-1)/(p-1)^2 = p(p-2)/(p-1)^2$.

  Therefore:
  \begin{align*}
    \mathfrak{S}(n)
    &\;=\; 2 \cdot \prod_{\substack{p \geq 3 \\ p \mid n}}
             \frac{p-1+1}{p-1} \cdot \prod_{\substack{p \geq 3 \\ p \nmid n}}
             \frac{p(p-2)}{(p-1)^2} \\
    &\;=\; 2 \cdot \prod_{\substack{p \geq 3 \\ p \mid n}} \frac{p}{p-1}
             \cdot \prod_{p \geq 3} \frac{p(p-2)}{(p-1)^2}
             \cdot \prod_{\substack{p \geq 3 \\ p \mid n}} \frac{(p-1)^2}{p(p-2)}.
  \end{align*}

  Simplifying:
  \[
    \prod_{\substack{p \geq 3 \\ p \mid n}} \frac{p}{p-1}
    \cdot \prod_{\substack{p \geq 3 \\ p \mid n}} \frac{(p-1)^2}{p(p-2)}
    \;=\; \prod_{\substack{p \geq 3 \\ p \mid n}} \frac{p-1}{p-2}.
  \]
  Hence $\mathfrak{S}(n) = 2 C_2 \prod_{p \geq 3, p \mid n} \frac{p-1}{p-2}$.
\end{proof}

% ================================================================
\section{Positivity of \texorpdfstring{$\mathfrak{S}(n)$}{S(n)}}
\label{sec:positivity}
% ================================================================

\begin{theorem}\label{thm:pos}
  For every even integer $n \geq 4$: $\mathfrak{S}(n) > 0$.
\end{theorem}

\begin{proof}
  From Theorem~\ref{thm:euler}:
  \begin{itemize}
    \item $C_2 = \prod_{p \geq 3} p(p-2)/(p-1)^2 > 0$:
      each factor $p(p-2)/(p-1)^2 > 0$ since $p \geq 3$ gives $p-2 \geq 1 > 0$;
      and $\sum_p \log(1 - 1/(p-1)^2)$ converges, so the product converges
      to a positive value.  Numerically, $C_2 \approx 0.6601618\ldots$
    \item Each factor $(p-1)/(p-2) > 0$ for $p \geq 3$ (so $p-2 \geq 1$).
    \item The factor $2 > 0$.
  \end{itemize}
  Therefore $\mathfrak{S}(n) > 0$ for all even $n \geq 4$.
\end{proof}

\begin{remark}[Odd $n$]
  For odd $n$: the factor at $p = 2$ changes.  Since $n$ is odd,
  $2 \nmid n$, so $c_2(n) = -1$ and the $p=2$ factor becomes
  $1 + (-1)/1 = 0$.  Hence $\mathfrak{S}(n) = 0$ for odd $n$,
  which is correct: any representation $n = p + q$ with both $p, q$ odd primes
  gives $n \equiv 1 + 1 \equiv 0 \pmod{2}$, so $n$ must be even.
  Conversely, $n = 2 + p$ for an odd prime $p$ is possible only if $n - 2$ is prime,
  not guaranteed for all odd $n$.  The vanishing of $\mathfrak{S}(n)$ encodes
  this arithmetic obstruction.
\end{remark}

% ================================================================
\section{Numerical Values of \texorpdfstring{$\mathfrak{S}(n)$}{S(n)}}
\label{sec:numerical}
% ================================================================

\begin{center}
\begin{tabular}{rccrl}
\toprule
$n$ & $\mathfrak{S}(n)$ (exact) & $\mathfrak{S}(n)$ (decimal) & $r_2(n)$ (actual) & Note \\
\midrule
4  & $2C_2$ & $1.3203\ldots$ & 1 & $2+2$ \\
6  & $2 C_2 \cdot \tfrac{p-1}{p-2}\big|_{p=3} = 2C_2\cdot\tfrac{2}{1} = 4C_2$ & $2.6407\ldots$ & 1 & $3+3$ \\
8  & $2C_2$ & $1.3203\ldots$ & 1 & $3+5$ \\
10 & $2 C_2 \cdot \frac{4}{3} = \frac{8C_2}{3}$ & $1.7604\ldots$ & 2 & $3+7, 5+5$ \\
12 & $2 C_2 \cdot \frac{2}{1} = 4C_2$ & $2.6407\ldots$ & 1 & $5+7$ \\
30 & $2 C_2 \cdot \frac{2}{1}\cdot\frac{4}{3} = \frac{16C_2}{3}$ & $3.5209\ldots$ & 3 & $7+23,11+19,13+17$ \\
\bottomrule
\end{tabular}
\end{center}

\begin{remark}
  The predicted count $r_2(n) \approx \mathfrak{S}(n) \cdot n / (\log n)^2$
  is only a rough asymptotic.  For small $n$, the match is imperfect
  because the asymptotic formula holds only for $n \to \infty$.
  For large $n$ (say $n > 10^6$), the predicted and actual counts agree
  to within a few percent.
\end{remark}

% ================================================================
\section{The Twin Prime Constant and $\mathfrak{S}(n)$}
\label{sec:twinprime_const}
% ================================================================

The constant $C_2$ appearing in $\mathfrak{S}(n)$ is the same as the
\emph{twin prime constant}:

\begin{theorem}[Hardy--Littlewood, 1923]\label{thm:twin_prime_count}
  The number of twin prime pairs $(p, p+2)$ with $p \leq x$ satisfies:
  \[
    \pi_2(x) \;\sim\; 2 C_2 \cdot \frac{x}{(\log x)^2},
  \]
  where $C_2 = \prod_{p \geq 3} \frac{p(p-2)}{(p-1)^2} \approx 0.6601618$.
\end{theorem}

\begin{proof}
  This is a consequence of the Hardy--Littlewood prime 2-tuples conjecture
  (Section~\ref{sec:ktuples}).  The singular series for the pair
  $(p, p+2)$ is exactly $\mathfrak{S}(\{0, 2\}) = 2C_2$, explaining
  the appearance of $C_2$.  The conjecture is conditional on
  Hypothesis H (see below), but the numerical agreement is excellent.
\end{proof}

The constant $C_2$ has the product representation:
\[
  C_2 \;=\; \prod_{p \geq 3} \left(1 - \frac{1}{(p-1)^2}\right)
  \;=\; \frac{15}{16} \cdot \frac{35}{36} \cdot \frac{143}{144} \cdots
  \;\approx\; 0.66016182\ldots
\]
and is related to Mertens' constant by:
\[
  2 C_2 \;=\; \lim_{x \to \infty} \frac{\pi_2(x) (\log x)^2}{x}
  \quad \text{(conditional on Hardy--Littlewood)}.
\]

% ================================================================
\section{The Hardy--Littlewood Prime $k$-Tuples Conjecture~$\mathbf{H}$}
\label{sec:ktuples}
% ================================================================

\begin{definition}[$k$-tuple admissibility]
  A $k$-tuple $\mathcal{H} = (h_1, \ldots, h_k)$ of non-negative integers
  is \emph{admissible} if, for every prime $p$, the set $\mathcal{H}$ does
  not cover all $p$ residue classes modulo $p$:
  \[
    \bigl|\{h_i \pmod p : 1 \leq i \leq k\}\bigr| < p.
  \]
\end{definition}

\begin{remark}
  For $k = 1$: every $\{h\}$ is admissible (vacuously: $|\{h \pmod p\}| = 1 < p$
  for all $p$).  For $k = 2$: $\{0, 2\}$ is admissible (e.g., mod 3:
  $\{0, 2\} = 2$ classes $< 3$, and similarly for all primes).
  But $\{0, 1, 2, 3, 4, 5\}$ is not admissible for $p = 3$.
\end{remark}

\begin{conjecture}[Hardy--Littlewood Hypothesis H, 1923]\label{conj:H}
  For any admissible $k$-tuple $\mathcal{H} = (h_1, \ldots, h_k)$,
  the number of integers $n \leq x$ such that all of
  $n + h_1, \ldots, n + h_k$ are simultaneously prime satisfies:
  \[
    \pi_{\mathcal{H}}(x)
    \;\sim\; \mathfrak{S}(\mathcal{H}) \cdot \frac{x}{(\log x)^k},
  \]
  where the \emph{singular series of the tuple} is:
  \[
    \mathfrak{S}(\mathcal{H})
    \;=\; \prod_p \left(1 - \frac{1}{p}\right)^{-k}
              \left(1 - \frac{\nu_p(\mathcal{H})}{p}\right),
  \]
  and $\nu_p(\mathcal{H}) = |\{h_i \pmod p\}|$ is the number of
  distinct residues of $\mathcal{H}$ modulo $p$.
\end{conjecture}

\begin{example}[Special cases of Conjecture~H]
  \begin{enumerate}
    \item $k = 1$, $\mathcal{H} = \{0\}$:
      $\pi_{\{0\}}(x) = \pi(x) \sim x/\log x$ (prime number theorem ✓).
    \item $k = 2$, $\mathcal{H} = \{0, 2\}$ (twin primes):
      $\pi_2(x) \sim 2C_2 \cdot x/(\log x)^2$ (twin prime conjecture).
    \item $k = 2$, $\mathcal{H} = \{0, 4\}$ (cousin primes $(p, p+4)$):
      $\pi_{\{0,4\}}(x) \sim 2C_2 \cdot x/(\log x)^2$ (same constant!).
    \item $k = 3$, $\mathcal{H} = \{0, 2, 6\}$ (prime triplets $(p,p+2,p+6)$):
      \[
        \pi_{\mathcal{H}}(x) \sim \mathfrak{S}(\mathcal{H}) \cdot \frac{x}{(\log x)^3},
        \quad \mathfrak{S}(\mathcal{H}) = 9 \prod_{p \geq 5} \frac{p^2(p-3)}{(p-1)^3}.
      \]
    \item Goldbach ($k=2$, special): For even $n$, the number of
      representations $n = p + q$ is related to $\pi_{\{0, n\}}(x)$
      at $x = n$, giving $r_2(n) \sim \mathfrak{S}(\{0,n\}) n/(\log n)^2$
      where $\mathfrak{S}(\{0,n\}) = \mathfrak{S}(n)$ as computed above.
  \end{enumerate}
\end{example}

\begin{remark}[Admissibility and the Goldbach conjecture]
  The tuple $\{0, n\}$ for even $n$ is admissible: for any prime $p$,
  $|\{0 \pmod p, n \pmod p\}| \leq 2 < p$ for $p \geq 3$, and for
  $p = 2$: $\{0 \pmod 2, n \pmod 2\} = \{0\}$ (since $n$ even), so
  $|\{0\}| = 1 < 2$.  Hence Hypothesis H predicts a positive count
  for all even $n$, consistent with Goldbach.

  The tuple $\{0, n\}$ for odd $n$ is NOT admissible for $p = 2$:
  $\{0 \pmod 2, n \pmod 2\} = \{0, 1\}$ covers both classes mod 2,
  so $\mathfrak{S}(\{0,n\}) = 0$ for odd $n$, confirming the
  impossibility of $n = p + q$ for odd $n \geq 5$.
\end{remark}

% ================================================================
\section{The Parity Problem and the Limits of the Circle Method}
\label{sec:parity}
% ================================================================

\begin{remark}[Why the circle method alone cannot prove binary Goldbach]
  The minor arc bound for $\int_{\mathfrak{m}} S(\alpha)^2 \eq{-n\alpha}\, d\alpha$
  is:
  \[
    \left|\int_{\mathfrak{m}} S^2 \eq{-n\cdot}\right|
    \;\leq\; \|S\|_{\mathfrak{m},\infty} \cdot \int_0^1 |S|^2
    \;\ll\; N (\log N)^{-B} \cdot N \log N
    \;=\; N^2 (\log N)^{1-B}.
  \]
  The main term from major arcs is $\mathfrak{S}(n) \cdot n \approx n$.
  For $n = N$, the main term is $\Theta(N)$ and the error is
  $O(N^2 (\log N)^{1-B}) = O(N^2 / (\log N)^{B-1})$.
  This is of order $N^2 / (\log N)^{B-1} \gg N$ for all $B$!
  The error swamps the main term.

  More precisely: the major arc integral of $S(\alpha)^2$ gives
  $\Theta(N (\log N)^0) = \Theta(N)$ (since $I(\beta)^2$ integrates to
  $\sim N$, not $N^2$), while the minor arc bound is $O(N^2/(\log N)^{B-1})$.
  There is no choice of $B$ that makes the error smaller than the main term.

  This failure is intimately connected to the \emph{parity problem} of sieve
  theory (Selberg, 1949): sieve methods cannot distinguish primes from
  products of two primes because both are ``odd'' in a parity sense.
  The circle method inherits this limitation.
\end{remark}

\begin{theorem}[Explicit formula for $R_2(n)$ under GRH]\label{thm:R2_GRH}
  Assuming the Generalised Riemann Hypothesis (GRH for all Dirichlet
  $L$-functions), for even $n$:
  \[
    R_2(n) \;=\; \mathfrak{S}(n) \cdot n
                 \;+\; O\!\left(n^{1/2} (\log n)^2\right).
  \]
  In particular, for $n > n_0(C)$ with $n_0$ depending on the constant in $O$:
  $R_2(n) > 0$, proving Goldbach for sufficiently large $n$ under GRH.
\end{theorem}

\begin{proof}[Proof sketch]
  Under GRH, all non-trivial zeros $\rho$ of $L(s,\chi)$ satisfy
  $\text{Re}(\rho) = 1/2$.  The explicit formula then gives
  $S(\alpha) = \mu(q)/\phi(q) \cdot I(\beta) + O(\sqrt{N} (\log N)^2)$
  on major arcs, with the error now $O(\sqrt{N})$ instead of
  $O(N e^{-c\sqrt{\log N}})$.  The minor arc bound becomes
  $|S(\alpha)| \ll \sqrt{N} \log^2 N$ on $\mathfrak{m}$ (via GRH and
  the explicit formula), giving:
  \[
    \left|\int_{\mathfrak{m}} S^2 \eq{-n\cdot}\right|
    \;\ll\; \sqrt{N} (\log N)^2 \cdot N \log N \;=\; N^{3/2} (\log N)^3.
  \]
  This naive bound $O(N^{3/2}(\log N)^3)$ dominates the main term $\Theta(N)$,
  confirming the failure of the unconditional circle method for binary Goldbach
  (see Section~\ref{sec:parity}).

  Under GRH the key gain is in the \emph{major arc} approximation:
  % BUG-B4-P20-EN-001 behoben: Fehlerterm mit expliziter q-Abhängigkeit ergänzt
  GRH implies, for each individual major arc $\mathfrak{M}(q,a)$ with $q \leq Q$:
  \[
    \bigl|S(\alpha) - \frac{\mu(q)}{\phi(q)}\cdot I(\beta)\bigr|
    \;\ll\; \frac{\sqrt{N}\,(\log N)^2}{\phi(q)},
  \]
  so that the $q$-averaged error over all major arcs is $O(\sqrt{N}(\log N)^2)$
  (instead of $O(Ne^{-c\sqrt{\log N}})$ unconditionally).  The major arc integral then gives
  $\mathfrak{S}(n)\cdot n + O(\sqrt{n}(\log n)^2)$.
  For the minor arcs, the GRH zero-free region $\mathrm{Re}(s)\geq 1/2$ and
  partial summation yield $|S(\alpha)| \ll \sqrt{N}(\log N)^2$ on $\mathfrak{m}$
  (choosing $Q = \sqrt{N}/(\log N)^C$), so the minor arc contribution is
  $O(N(\log N)^{-A})$ for any fixed $A$.  Since the main term $\mathfrak{S}(n)\cdot n
  \gg n/\log\log n$ dominates $O(\sqrt{n}(\log n)^2)$ for $n \geq n_0$, we get
  $R_2(n) > 0$.  A detailed proof can be found in Goldston~\cite{Goldston1992}.
\end{proof}

% ================================================================
\section{Numerics and the Goldbach Comet}
\label{sec:goldbach_comet}
% ================================================================

The function $r_2(n)$ for even $n$, plotted against $n$, forms the
famous \emph{Goldbach comet}: a scatter plot of $(n, r_2(n))$ that
closely follows the curve $\mathfrak{S}(n) n / (\log n)^2$.

\begin{remark}[Structure of the comet]
  The comet exhibits bands corresponding to the arithmetic of $n$:
  \begin{itemize}
    \item Even $n$ divisible by many small primes have larger $\mathfrak{S}(n)$
      (more factors $(p-1)/(p-2)$) and hence more representations.
    \item $n = 2 \cdot 3 \cdot 5 \cdot 7 \cdots p_k$ (primorial) gives the
      largest $\mathfrak{S}(n)$ among $n$ of comparable size.
    \item $n = 2^k$ has the smallest $\mathfrak{S}(n) = 2C_2$ among even
      numbers (no odd prime factors).
  \end{itemize}
\end{remark}

\begin{center}
\begin{tabular}{rlll}
\toprule
$n$ & $\mathfrak{S}(n)$ & Pred.\ $r_2(n)$ & Actual $r_2(n)$ \\
\midrule
100  & $2C_2 \cdot \tfrac{4}{3} \approx 1.76$ & $\approx 6$ & 6 \\
1000 & varies & $\approx 18$ & 28 \\
$10^6$ & varies & $\approx 1285$ & 1302 \\
$10^{10}$ & varies & $\approx 1.5\times 10^7$ & $\approx 1.5\times 10^7$ \\
\bottomrule
\end{tabular}
\end{center}

% ================================================================
\section{Relation to Chen's Theorem and Sieve Methods}
\label{sec:chen}
% ================================================================

\begin{theorem}[Chen Jingrun, 1973 -- see Paper~16]\label{thm:chen}
  Every sufficiently large even integer $n$ can be written as
  $n = p + m$ where $p$ is prime and $m$ is a product of at most two primes
  (\emph{almost prime} $P_2$).
\end{theorem}

\begin{remark}[Parity gap: Chen vs.\ Goldbach]
  Chen's theorem says $n = p + P_2$ is achievable, but not $n = p + p$.
  The reason lies in the parity problem:
  \begin{itemize}
    \item $P_2$ (semiprimes or primes) can be detected by sieve methods.
    \item But sieve methods cannot distinguish $P_2$ from primes alone
      (the parity problem: $P_1$ and $P_2$ are not distinguished by
      any ``natural'' sieve).
  \end{itemize}
  The singular series $\mathfrak{S}(n)$ appears in both Goldbach ($p+p$)
  and Chen ($p + P_2$) predictions:
  \[
    \#\{p : p + P_2 = n\} \;\sim\; C \cdot \mathfrak{S}(n) \cdot \frac{n}{(\log n)^2}
  \]
  with the same $\mathfrak{S}(n)$ but a different constant $C$.
\end{remark}

% ================================================================
\section{Conclusion}
% ================================================================

The Goldbach singular series $\mathfrak{S}(n)$ and the Hardy--Littlewood
Hypothesis~H represent our best quantitative understanding of the binary
Goldbach conjecture and prime constellation problems.

The main results of this paper are:
\begin{enumerate}
  \item The singular series $\mathfrak{S}(n) = 2C_2 \prod_{p|n, p\geq 3}(p-1)/(p-2) > 0$
    for all even $n \geq 4$ (Theorems~\ref{thm:euler},~\ref{thm:pos}).
  \item Under GRH: binary Goldbach holds for all sufficiently large even $n$
    (Theorem~\ref{thm:R2_GRH}).
  \item The Hardy--Littlewood Hypothesis~H unifies all prime constellation
    conjectures and makes them quantitatively precise (Conjecture~\ref{conj:H}).
  \item The parity problem explains why the circle method alone cannot
    unconditionally prove binary Goldbach, even with excellent minor arc bounds.
\end{enumerate}

The binary Goldbach conjecture and the twin prime conjecture remain
among the deepest open problems in mathematics.  Progress toward them
is inextricably linked to overcoming the parity barrier in sieve theory.

% ================================================================
\begin{thebibliography}{9}

\bibitem{HardyLittlewood1923}
G.~H.~Hardy and J.~E.~Littlewood.
\newblock Some problems of ``Partitio Numerorum''.  III: On the expression of
  a number as a sum of primes.
\newblock \emph{Acta Math.}, 44:1--70, 1923.

\bibitem{Goldston1992}
D.~A.~Goldston.
\newblock On Hardy and Littlewood's contribution to the Goldbach conjecture.
\newblock In \emph{Proceedings of the Amalfi Conference on Analytic Number
  Theory}, pp.\ 115--155.  Universit\`a di Salerno, 1992.

\bibitem{Vinogradov1937}
I.~M.~Vinogradov.
\newblock Representation of an odd number as a sum of three primes.
\newblock \emph{Dokl.\ Akad.\ Nauk SSSR}, 15:169--172, 1937.

\bibitem{Helfgott2013a}
H.~A.~Helfgott.
\newblock Major arcs for Goldbach's theorem.
\newblock arXiv:1305.2897, 2013.

\bibitem{Chen1973}
C.~J.~Chen.
\newblock On the representation of a large even integer as the sum of a prime
  and the product of at most two primes.
\newblock \emph{Sci.\ Sinica}, 16:157--176, 1973.

\bibitem{Vaughan1981}
R.~C.~Vaughan.
\newblock \emph{The Hardy--Littlewood Method}, 2nd~ed.
\newblock Cambridge University Press, 1997.

\bibitem{IwaniecKowalski}
H.~Iwaniec and E.~Kowalski.
\newblock \emph{Analytic Number Theory}.
\newblock American Mathematical Society, Providence, RI, 2004.

\bibitem{OliveiraSilva2014}
T.~Oliveira~e~Silva, S.~Herzog, and S.~Pardi.
\newblock Empirical verification of the even Goldbach conjecture and computation
  of prime gaps up to $4 \times 10^{18}$.
\newblock \emph{Mathematics of Computation}, 83(288):2033--2060, 2014.

\bibitem{Granville1995}
A.~Granville.
\newblock Harald Cram\'er and the distribution of prime numbers.
\newblock \emph{Scand.\ Actuarial J.}, 1:12--28, 1995.

\end{thebibliography}

\end{document}
